% Options for packages loaded elsewhere
\PassOptionsToPackage{unicode}{hyperref}
\PassOptionsToPackage{hyphens}{url}
%
\documentclass[
  10pt,
  english,
]{article}
\usepackage{amsmath,amssymb}
\usepackage{lmodern}
\usepackage{iftex}
\ifPDFTeX
  \usepackage[T1]{fontenc}
  \usepackage[utf8]{inputenc}
  \usepackage{textcomp} % provide euro and other symbols
\else % if luatex or xetex
  \usepackage{unicode-math}
  \defaultfontfeatures{Scale=MatchLowercase}
  \defaultfontfeatures[\rmfamily]{Ligatures=TeX,Scale=1}
\fi
% Use upquote if available, for straight quotes in verbatim environments
\IfFileExists{upquote.sty}{\usepackage{upquote}}{}
\IfFileExists{microtype.sty}{% use microtype if available
  \usepackage[]{microtype}
  \UseMicrotypeSet[protrusion]{basicmath} % disable protrusion for tt fonts
}{}
\makeatletter
\@ifundefined{KOMAClassName}{% if non-KOMA class
  \IfFileExists{parskip.sty}{%
    \usepackage{parskip}
  }{% else
    \setlength{\parindent}{0pt}
    \setlength{\parskip}{6pt plus 2pt minus 1pt}}
}{% if KOMA class
  \KOMAoptions{parskip=half}}
\makeatother
\usepackage{xcolor}
\usepackage[skins,breakable]{tcolorbox}
\usepackage{titlesec}
\definecolor{paperink}{HTML}{26231F}
\definecolor{paperaccent}{HTML}{6E5943}
\definecolor{paperline}{HTML}{C8BCAD}
\definecolor{paperboxbg}{HTML}{F7F3EC}
\color{paperink}
\titleformat{\subsection}{\normalfont\large\bfseries\color{paperink}}{}{0pt}{\color{paperaccent}\rule{2pt}{1.05em}\hspace{7pt}}
\titlespacing*{\subsection}{0pt}{10pt}{5pt}
\tcbset{paperbox/.style={enhanced,breakable,colback=paperboxbg,colframe=paperline,
  boxrule=.45pt,arc=2mm,left=8pt,right=8pt,top=6pt,bottom=6pt,
  before skip=6pt,after skip=8pt},
  accentbox/.style={paperbox,borderline west={2pt}{0pt}{paperaccent}},
  compactbox/.style={paperbox,top=3pt,bottom=3pt,before skip=3pt,after skip=5pt},
  compactaccent/.style={compactbox,borderline west={2pt}{0pt}{paperaccent}}}
\IfFileExists{xurl.sty}{\usepackage{xurl}}{} % add URL line breaks if available
\IfFileExists{bookmark.sty}{\usepackage{bookmark}}{\usepackage{hyperref}}
\hypersetup{
  pdftitle={Certified Coverage of Exponential Search Spaces},
  pdfauthor={J. R. Landers},
  pdflang={en},
  hidelinks,
  pdfcreator={LaTeX via pandoc}}
\urlstyle{same} % disable monospaced font for URLs
\usepackage[margin=0.72in]{geometry}
\setlength{\emergencystretch}{3em} % prevent overfull lines
\providecommand{\tightlist}{%
  \setlength{\itemsep}{0pt}\setlength{\parskip}{0pt}}
\setcounter{secnumdepth}{-\maxdimen} % remove section numbering
\ifXeTeX
  % Load polyglossia as late as possible: uses bidi with RTL langages (e.g. Hebrew, Arabic)
  \usepackage{polyglossia}
  \setmainlanguage[]{english}
\else
  \usepackage[main=english]{babel}
% get rid of language-specific shorthands (see #6817):
\let\LanguageShortHands\languageshorthands
\def\languageshorthands#1{}
\fi
\ifLuaTeX
  \usepackage{selnolig}  % disable illegal ligatures
\fi

\title{Certified Coverage of Exponential Search Spaces}
\author{J. R. Landers}
\date{September 2026}

\begin{document}
\maketitle
\vspace{-1.4em}
\noindent\textcolor{paperaccent}{\rule{.11\linewidth}{1.5pt}}\textcolor{paperline}{\rule{.89\linewidth}{.45pt}}
\vspace{.4em}

\begin{tcolorbox}[accentbox]
\noindent\textbf{Idea.} Consider a search problem whose naive state space has size
\(2^n\). Some states may admit a compact structural representation;
others may be proposed for elimination by a learned model and certified
exactly; still others may be represented by exact summaries of paired
groups. Only the residual must be searched explicitly. The central
quantity is the \emph{growth of what remains uncovered}, together with
the cost of obtaining coverage.
\end{tcolorbox}

\hypertarget{a-coverage-model}{%
\subsection{A coverage model}\label{a-coverage-model}}

Let \(\mathcal{S}_n\) be a nominal search space with

\[
|\mathcal{S}_n| = 2^n .
\]

We use three mechanisms to cover disjoint parts of this space.

First, a structural procedure identifies a set
\(K_n \subseteq \mathcal{S}_n\) of size

\[
k(n) = |K_n|
\]

that admits a compact representation. In the simplest logarithmic model,
the \(k(n)\) explicit states are replaced by a description of cost

\[
C_{\mathrm{comp}}(n) = \log_2 k(n).
\]

This is deliberately an abstract assumption: it describes the cost once
such a representation is available, not a universal property of
arbitrary subsets.

Second, a learned classifier proposes regions of the remaining space for
removal. Let \(\delta(n)\) denote the total number of states covered by
these proposals. The classifier is not trusted to delete states
directly. Instead, an exact verifier certifies proposed regions. If
\(a(n) \in [0,1]\) is the fraction of proposed mass that is successfully
certified, then the safely eliminated mass is

\[
\delta_c(n) = a(n)\,\delta(n).
\]

Prediction error therefore affects efficiency, not correctness:
uncertified proposals simply remain in the residual problem.

Third, pair groups of states and replace each pair by an exact summary
for the search question. Let \(p(n)\) count the additional states
covered by such summaries, excluding states already counted in \(k(n)\)
or \(a(n)\delta(n)\). Let \(Q_{\mathrm{pair}}(n)\) include the cost of
finding pairs, constructing and certifying their summaries, and using
the resulting representation.

The explicit residual is

\begin{tcolorbox}[accentbox]
\begin{equation}
r(n) = 2^n - k(n) - a(n)\delta(n) - p(n).
\end{equation}
\end{tcolorbox}

Let \(P(n)\) collect the cost of discovering the structural
representation, running the learned proposal mechanism, certifying its
proposed regions, and managing the resulting representation. Under the
simple assumption that explicit residual search costs \(r(n)\) units, the
total modeled cost is

\begin{tcolorbox}[paperbox]
\begin{equation}
T(n) = P(n) + Q_{\mathrm{pair}}(n) + r(n) + \log_2 k(n).
\end{equation}
\end{tcolorbox}

Thus

\begin{equation}
\begin{aligned}
T(n) ={}& P(n) + Q_{\mathrm{pair}}(n) + 2^n - k(n)\\
      &{}- a(n)\delta(n) - p(n) + \log_2 k(n).
\end{aligned}
\end{equation}

The separation is important. Large covered regions help only if
discovering, representing, and certifying them is cheap.

\hypertarget{the-certified-coverage-theorem}{%
\subsection{The certified coverage
theorem}\label{the-certified-coverage-theorem}}

Define the total \emph{certified coverage}

\[
C(n) = k(n) + a(n)\delta(n) + p(n),
\]

so that \(r(n) = 2^n - C(n)\).

\begin{tcolorbox}[accentbox]
\textbf{Theorem 1 (Certified coverage criterion).} Assume
\(P(n)+Q_{\mathrm{pair}}(n) = \mathrm{poly}(n)\),
\(1 \le k(n) \le 2^n\), and the residual is searched in time
proportional to \(r(n)\). Then

\[
  T(n) = \mathrm{poly}(n) \quad\Longleftrightarrow\quad 2^n - C(n) = \mathrm{poly}(n).
  \]

Equivalently, polynomial modeled runtime occurs exactly when

\[
  C(n) = 2^n - \mathrm{poly}(n).
  \]
\end{tcolorbox}

{Proof.} From (2),

\[
  T(n) = P(n) + Q_{\mathrm{pair}}(n) + \bigl(2^n - C(n)\bigr) + \log_2 k(n).
  \]

Because \(k(n) \le 2^n\),

\[
  0 \le \log_2 k(n) \le n.
  \]

The overhead and logarithmic terms are polynomially bounded. Hence
\(T(n)\) is polynomial exactly when the residual \(2^n - C(n)\) is
polynomially bounded. \hfill $\square$

The theorem gives a sharp asymptotic warning: eliminating a very large
\emph{percentage} of an exponential space need not change its asymptotic
order.

{Corollary 1 (Fixed uncovered fraction).} If there exists a constant
\(\varepsilon > 0\) such that, for infinitely many \(n\),

\[
  C(n) \le (1-\varepsilon)2^n,
  \]

then

\[
  r(n) \ge \varepsilon 2^n
  \]

and therefore, along that infinite subsequence, the modeled runtime satisfies

\[
  T(n) = \Omega(2^n)
  \]

Thus 90\%, 99\%, or even 99.999\% coverage is asymptotically
insufficient if the uncovered fraction stays bounded away from zero.
Polynomial behavior requires

\[
\frac{C(n)}{2^n} = 1 - \frac{\mathrm{poly}(n)}{2^n},
\]

so the uncovered fraction must vanish exponentially fast, up to
polynomial factors.

\hypertarget{differential-form}{%
\subsection{Differential form}\label{differential-form}}

The discrete complexity statement has a useful continuous interpolation.
Treat \(n\) as a continuous scale parameter and differentiate (1):

\begin{tcolorbox}[compactbox]
\begin{equation}
\frac{dr}{dn} = (\ln 2)2^n - k'(n) - a(n)\delta'(n) - a'(n)\delta(n) - p'(n).
\end{equation}
\end{tcolorbox}

The terms are raw state creation, structural coverage, new certified
proposals, improving certification rate, and paired coverage,
respectively.

Differentiating the full cost gives

\begin{tcolorbox}[compactaccent]
\begin{equation}
\begin{aligned}
\frac{dT}{dn} ={}& (\ln 2)2^n - \left(1 - \frac{1}{k \ln 2}\right)k' - a\delta' - a'\delta - p'\\
&{}+ P'(n) + Q_{\mathrm{pair}}'(n).
\end{aligned}
\end{equation}
\end{tcolorbox}

For large \(k\), the \(1/(k\ln 2)\) correction is negligible, giving the
approximate balance law

\begin{tcolorbox}[compactbox]
\begin{equation}
\begin{aligned}
T'(n) \approx{}& (\ln 2)2^n - k' - a\delta' - a'\delta - p'\\
&{}+ P'(n) + Q_{\mathrm{pair}}'(n).
\end{aligned}
\end{equation}
\end{tcolorbox}

Equivalently, since \(r(n) = 2^n - C(n)\),

\[
r'(n) = (\ln 2)2^n - C'(n).
\]

Integrating from a reference scale \(n_0\),

\begin{tcolorbox}[compactbox]
\begin{equation}
r(n) = r(n_0) + \int_{n_0}^{n} \bigl[(\ln 2)2^s - C'(s)\bigr]\,ds.
\end{equation}
\end{tcolorbox}

This form makes the central picture especially transparent: polynomial
behavior requires the cumulative gap between exponential state creation
and certified coverage growth to remain only polynomial in \(n\).

\hypertarget{why-certification-matters}{%
\subsection{Why certification matters}\label{why-certification-matters}}

A naive learned-pruning model might discard \(\delta(n)\) states with
error rate \(e(n)\). If every erroneous deletion has to be recovered to
preserve exactness, an error residue of order \(e(n)\delta(n)\) appears.
When \(\delta(n) = \Theta(2^n)\), any fixed \(e > 0\) leaves an
exponential residue. To keep that residue polynomial would require

\[
e(n)\,2^n = \mathrm{poly}(n), \qquad\text{hence}\qquad e(n) \le \frac{\mathrm{poly}(n)}{2^n}.
\]

Thus ordinary constant-error prediction is not enough for exact bulk
deletion of an exponential region.

Certification avoids that failure mode. The learned model becomes a
\emph{proposal mechanism}: it suggests large regions that may be
irrelevant, while an exact verifier alone has authority to delete them.
A poor classifier can waste computation, but it cannot make the final
algorithm incorrect. In this formulation, learning changes \(\delta\)
and the acceptance function \(a\); correctness is carried by
verification.

\hypertarget{the-cost-of-pairing}{%
\subsection{The cost of pairing}\label{the-cost-of-pairing}}

For disjoint groups \(B_0,B_1\), a summary \(s\) must support an exact
merge rule

\begin{tcolorbox}[accentbox]
\[
s(B_0\cup B_1)=\operatorname{Merge}\bigl(s(B_0),s(B_1)\bigr),
\]
\end{tcolorbox}

and the final summary must answer the original search question exactly.
A complete explicit pairing tree over \(2^n\) states has \(2^n-1\)
merges, even though only one group remains. Thus pairing helps
asymptotically only when a compact rule or reuse of summaries keeps
\(Q_{\mathrm{pair}}(n)\) polynomial. Counting covered states without
this cost would hide the original exponential work.

\hypertarget{scope-and-the-remaining-difficulty}{%
\subsection{Scope and the remaining
difficulty}\label{scope-and-the-remaining-difficulty}}

The coverage criterion is a theorem about the model; it does not assert
that every concrete search problem admits such coverage. The hard
question is whether, for a given problem, one can simultaneously achieve

\[
P(n)+Q_{\mathrm{pair}}(n) = \mathrm{poly}(n) \qquad\text{and}\qquad 2^n - C(n) = \mathrm{poly}(n).
\]

Cardinality alone cannot answer this: exponentially many states can
sometimes be described by a short rule. The substantive lower- or
upper-bound problem is therefore about the \emph{cost of discovering and
certifying structure}, not merely the number of states represented by
that structure.

\hypertarget{general-representation-cost}{%
\subsection{General representation
cost}\label{general-representation-cost}}

The logarithmic term is the motivating special case, not an essential
part of the coverage theorem. If a compressed region of size \(k(n)\)
has representation cost \(\rho(n,k)\), write

\[
T(n) = P(n) + Q_{\mathrm{pair}}(n) + r(n) + \rho(n,k(n)).
\]

Whenever \(P(n)+Q_{\mathrm{pair}}(n)\) and \(\rho(n,k(n))\) are
polynomially bounded, the same proof gives

\[
T(n) = \mathrm{poly}(n) \quad\Longleftrightarrow\quad r(n) = 2^n - C(n) = \mathrm{poly}(n).
\]

Thus the essential object is the uncovered residual; \(\log_2 k\) is
simply the strongest compression regime considered here. This also
separates two questions that should not be conflated: how many states a
representation covers, and how costly that representation is to discover
and manipulate.

\begin{tcolorbox}[accentbox]
\textbf{Summary.} Exponential search becomes a competition between state-space
growth and certified coverage. Structural compression and exact pairing
represent states collectively; learning proposes further regions;
certification makes eliminations exact; and the residual determines
explicit work. The governing condition is simple:

\[
  \text{polynomial overhead} + \text{polynomial residual} \;\Longrightarrow\; \text{polynomial modeled runtime.}
  \]

The differential form records the same statement dynamically: coverage
must track exponential state creation closely enough that their
accumulated gap is only polynomial.
\end{tcolorbox}

\end{document}
